\documentclass[11pt]{article}
\usepackage{fontspec}
\usepackage[french,english]{babel}
\usepackage{amsmath}
\usepackage{amssymb}
\usepackage{amsthm}
\usepackage{mathrsfs}
\usepackage{tikz-cd}
\usepackage{enumitem}

\newcommand{\printedpage}[1]{\marginpar{\footnotesize p.~#1}}
\newcommand{\msp}[1]{\marginpar{\footnotesize [#1]}}
\newcommand{\Ob}{\operatorname{Ob}}
\newcommand{\Hom}{\operatorname{Hom}}
\newcommand{\Ind}{\operatorname{Ind}}
\newcommand{\card}{\operatorname{card}}
\newcommand{\cM}{\mathcal{M}}
\newcommand{\cE}{\mathcal{E}}
\newcommand{\cC}{\mathcal{C}}
\newcommand{\cQ}{\mathcal{Q}}
\newcommand{\uS}{\underline{S}}
\newcommand{\uT}{\underline{T}}
\newcommand{\U}{\mathfrak{U}}

\begin{document}

\setcounter{footnote}{195}

\printedpage{184}
% head: A. Grothendieck : Dérivateurs, Ch. XVIII, Édition des Universités de Montpellier II et Paris VII
% foot: Version 17 décembre 2025

\msp{page 475}
We then find the

\begingroup\itshape
\noindent\textbf{\upshape Corollary 4.16.6.2.} Let $f : \cE \longrightarrow \cC$ be a functor. Equivalent conditions:
\begin{enumerate}[label={}]
\item[a)] For every category $\cC'$ that is \textup{small} and \textup{Karoubian}, and every functor $\cC' \longrightarrow \cC$, $\cE' = \cE \times_{\cC} \cC'$ is accessible, and the functor $\cE' \longrightarrow \cC'$ is accessible.
\item[b)] As in a), but $\cC'$ being one of the three categories $\Delta_1$, $\cQ$, $\cQ'$ (where $\cQ$, $\cQ'$ are defined in 4.15.6.1 above).
\item[b$'$)] \textup{(}\footnote{?, to be checked carefully!}\textup{)}. (If $\cC$ is without projectors.) As in b), but restricting to the case $\cC' = \boldsymbol{\Delta}_1$ (to the exclusion of $\cQ$, $\cQ'$).
\end{enumerate}
\endgroup

\begingroup\itshape
\noindent\textbf{\upshape Corollary 4.16.6.3.} Suppose $f : \cE \longrightarrow \cC$ is fibrant or cofibrant. Then the preceding equivalent conditions imply that:
\begin{enumerate}[label={}]
\item[1$^\circ$)] The $\cE_s$ are accessible ($s \in \Ob \cC$).
\item[2$^\circ$)] For every arrow $\lambda : s \longrightarrow t$ in $\cC$, the functor $\lambda^* : \cE_t \longrightarrow \cE_s$ (resp. $\cE_s \longrightarrow \cE_t$) is accessible.
\end{enumerate}
\noindent On the other hand, the collection of these conditions 1$^\circ$, 2$^\circ$) is equivalent to the validity of b) above, in the case $\cC' = \boldsymbol{\Delta}_1$, or again to the validity
\msp{page 476}
of a) above, when $\cC'$ is \textup{without projectors}.
\endgroup

\textbf{Question 4.16.6.4.} \textit{Do the preceding conditions 1$^\circ$ and 2$^\circ$ imply those of 4.16.6.2? In other words, if the base category $\cC$ is equal to $\cQ$ or to $\cQ'$, then do these conditions imply that $\cE$ is accessible, and $f : \cE \longrightarrow \cC$ accessible?}

It seems stupid not to know the answer!

\textbf{4.16.6.5.} I feel like calling a category \textit{accessible over $\cC$}, a category $\cE$ over $\cC$ satisfying the equivalent conditions of 4.16.6.1 -- at least in the case where $\cC$ is assumed small. For this terminology to be really satisfactory, we would have to know how to prove that these conditions imply this: for every accessible category $\cM$ and every accessible functor $\cM \longrightarrow \cC$, the category $\cE \times_{\cC} \cM$ is accessible, and the functor $\cE \times_{\cC} \cM \longrightarrow \cM$ is accessible. % [sic] satisfaisant -> satisfaisante, accord avec « terminologie »

\msp{page 476bis}
The idea is that one would call (under a smallness hypothesis) \textit{accessible category $\cE$} over another $\cC$, a category having the property we have just stated -- but now assuming in addition that $\cC$ satisfies $\mathrm{L}_\pi$, I imagine, for every accessible functor $\cM \longrightarrow \cC$, with $\cM$ accessible, $\cE \times_{\cC} \cM$ is accessible and the canonical functor $\cE \times_{\cC} \cM \longrightarrow \cE$ is accessible.

\msp{page 477}
\subsection*{4.16.7. Return to the closed subcategories of an accessible category (4.16.1).}

\textbf{Proposition 4.16.7.1.} \textit{Let $\cM$ be an accessible category, $(\uS_i)_{i\in I}$ a family of closed subcategories accessible in $\cM$. We do not assume that \textup{card}\,$I$ is small. Let $\uS$ be the strictly full subcategory of $\cM$ such that $\Ob \uS = \bigcup_{i\in I} \Ob \uS_i$. Then for $\uS$ % [sic] On ne suppose que -> On ne suppose pas que, « pas » apparemment omis
\printedpage{185}
% head: A. Grothendieck : Dérivateurs, Ch. XVIII, Édition des Universités de Montpellier II et Paris VII
% foot: Version 17 décembre 2025
to be accessible in $\cM$, it is necessary and sufficient that there exist a \textup{small} subset $J \subset I$, such that $\Ob \uS = \bigcup_{i\in J} \uS_i$. (In particular, if $I$ is ordered, and $i \longmapsto \uS_i$ increasing and large relative to every small cardinal, it is necessary and sufficient that there exist an $i \in I$ such that $\uS = \uS_i$, i.e. such that $j \geqslant i$ implies $\uS_j = \uS_i$.)}

This is obviously sufficient, by virtue of 4.13\dots{} Let us prove conversely that if $\uS$ is accessible in $\cM$, then the preceding condition is satisfied. Let $\pi$ be a cardinal $\geqslant \pi_0$ such that $\uS$ is $\pi$-accessible in $\cM$, hence
\[
\uS \ \xleftarrow{\ \sim\ }\ \Ind_\pi \uS(\pi) \ , \qquad\qquad \text{where } \uS(\pi) = \uS \cap \cM_\pi.
\]

\msp{page 478}
Let $\uS(\pi)_0$ be a skeletal subcategory of $\uS(\pi)$, such that $\uS(\pi)_0 \hookrightarrow \uS(\pi)$ is an equivalence, so that $\uS(\pi)_0$ is small. For every $x \in \Ob \uS(\pi)_0$, let $i_x \in I$ be such that $x \in \Ob \uS_{i_x}$. Let $J$ be the set of all the $i_x$, which is a small subset of $I$, and let $\uS_J$ be the strictly full subcategory of $\cM$ with set of objects $\bigcup_{j\in J} \Ob \uS_j$. Thus $\uS(\pi)_0$ is contained in $\uS_J$, hence also $\uS(\pi)$, since $\uS_J$ is strictly full: % [sic] que est -> qui est, pronom relatif sujet
\[
\uS(\pi) \ \hookrightarrow\ \uS_J \ .
\]
Moreover $\uS_J$ is accessible in $\cM$ (loc. cit.), and is closed in $\cM$ as a union of closed subcategories, \textit{a fortiori it satisfies $L_\pi$}. Thus, since $\uS(\pi) \subset \uS_J$, the essentially surjective functor $\Ind_\pi(\uS(\pi)) \longrightarrow \uS$ factors through $\uS_J$. Hence $\uS_J \hookrightarrow \uS$ is essentially surjective, and since $\uS_J$ is strictly full, we conclude that $\uS_J = \uS$, q.e.d.

\msp{page 479}
\textbf{Proposition 4.16.7.2.} \textit{Let $\cM$ be an accessible category, $x$ an object of $\cM$, and ${}_x\cM$ the closed subcategory of $\cM$ generated by $x$, consisting of the $y \in \Ob \cM$ such that $\Hom(x,y) \neq \emptyset$, i.e. the essential image of $x\backslash\cM \longrightarrow \cM$. Then ${}_x\cM$ is a subcategory of $\cM$ accessible in $\cM$.}

Indeed, it is the inverse image of the subcategory $\mathrm{Ens}^*$ accessible in $\mathrm{Ens}$ under the representable (hence accessible) functor
\[
h_x \ :\ y \longmapsto \Hom(x,y) \ :\ \cM \ \longrightarrow\ \mathrm{Ens} \ ,
\]
hence is accessible in $\cM$ by virtue of 4.13\dots

\textbf{Corollary 4.16.7.3.} \textit{Let $\cM$ be an accessible category, $\uS$ a closed subcategory. For $\uS$ to be accessible, it is necessary and sufficient that there exist a \textup{small} subset $J$ of $\uS$, such that $\Ob \uS = \bigcup_{x\in J} \Ob\,(\,{}_x\cM)$.}

Indeed, we have
\[
\Ob \uS \ =\ \bigcup_{x\in\Ob\uS} \Ob\,(\,{}_x\cM) \ ,
\]
i.e. $\uS$ appears as the union of
\msp{page 480}
the family (not necessarily small) of closed subcategories ${}_x\cM$ of $\cM$ (indexed by $x \in \Ob \uS$), accessible in $\cM$ by virtue of 4.16.7.2. Thus our assertion is a particular case of 4.16.7.1.

\printedpage{186}
% head: A. Grothendieck : Dérivateurs, Ch. XVIII, Édition des Universités de Montpellier II et Paris VII
% foot: Version 17 décembre 2025

One can also, for every small subset $J$ of $\Ob \uS$, consider the closed subcategory $\uS_J \subset \uS$ of $\cM$ generated by $J$. The set $I$ of the subsets $J$ in question, ordered by inclusion, is large relative to every cardinal $\pi \in \U$, and $\uS$ is an increasing union of the $\uS_J$. On the other hand, the proposition tells us that $\uS$ is accessible in $\cM$ if and only if there exists a $J \in I$, such that $\uS_J = \uS$, i.e. such that $\uS_J = \uS_{J'}$ for $J \subset J'$.

\begingroup\itshape
\noindent\textbf{\upshape Proposition 4.16.7.4.} Let $\cM$ be an accessible category. The following conditions are equivalent:
\begin{enumerate}[label={}]
\item[a)] There exists a subcategory $\uS$ of $\cM$ closed in $\cM$, and not accessible in $\cM$.
\item[b)] Let $I$ be the set of ordinals $\alpha$ such that $\card \alpha \in \U$, i.e. of the
\msp{page 481}
ordinals $\alpha$ which are in $\U$. Thus $I = I_{\alpha_0}$, where $\alpha_0$ is the smallest ordinal which does not belong to $\U$, and $I$ is large relative to every small cardinal $\pi$, i.e. relative to every cardinal in $\U$ \textup{(}\footnote{introduce $I = I_{\U}$ before the statement of the proposition.}\textup{)}. The condition under consideration is the following:

There exists a transfinite sequence % [sic] transfini -> transfinie, accord avec « suite »
\[
\underline{x} \ =\ (x_\alpha)_{\alpha\in I}
\]
of elements of $\Ob \cM$, such that we have
\[
\Hom(x_\beta, x_\alpha) \ =\ \emptyset \qquad\qquad \textit{if } \beta < \alpha.
\]
To be precise, for such a transfinite sequence, let $\uS$ or ${}_{\underline{x}}\cM$ be the closed subcategory of $\cM$ generated by the family of the $x_\alpha$, i.e. such that $\Ob \uS = \bigcup_{\alpha\in I} \Ob\ {}_{x_\alpha}\cM = \bigcup_{\alpha\in I} x_\alpha\backslash\Ob \cM$. Then ${}_{\underline{x}}\cM$ is a closed subcategory not accessible in $\cM$, and every closed subcategory not accessible in $\cM$ is obtained in this way.
\end{enumerate}
\endgroup

\textsc{Proof}. Let us consider a sequence $\underline{x} = (x_\alpha)_{\alpha\in I}$ as above. If $J$ is a small subset of $I$, there exists an $\alpha \in I$ which is an upper bound for the $\beta \in J$, hence $\Hom(x_\beta, x_\alpha) = \emptyset$, for $\beta \in J$,
\msp{page 482}
which means that $x_\alpha$ does not belong to any of the $\uS_\beta = {}_{x_\beta}\cM$, for $\beta \in J$. Thus $\uS$ is not the union of a small subfamily of the family of the ${}_{x_\alpha}\cM$, hence by virtue of 4.16.7.1, $\uS$ is not accessible in $\cM$. % [sic] n'est par réunion -> n'est pas réunion ; n'est par accessible -> n'est pas accessible, « par » pour « pas » (2 fois)

Conversely, let $\uS$ be such a subcategory in $\cM$. Let us consider a well-order on $\Ob \uS$ (objects $s_\alpha$), and construct the sequence of the $x_\alpha$ ($\geqslant s_\alpha$) by transfinite induction (\footnote{\textbf{N.B.} We have $\card \uS \notin \U$, otherwise, $\uS$ would be small, hence accessible in $\cM$, contrary to the hypothesis.}).

1) We choose $x_0 =$ first element of $\Ob \uS$ (exists, since $\Ob \uS \neq \emptyset$, i.e. $\uS \neq \emptyset$, since otherwise, $\uS$ would be accessible in $\cM$).

2) If the $x_\beta$ ($\beta < \alpha$) have been constructed, we take for $x_\alpha$ the first element of $\uS$ bounding the $x_\beta$ and which does not belong to the closed subcategory of $\cM$ generated by the $x_\beta$ ($\beta < \alpha$). (This closed subcategory is $\subset \uS$, and $\neq \uS$, since $\uS$ is not accessible in $\cM$.) Thus we do have
\[
\Hom(x_\beta, x_\alpha) \ =\ \emptyset \qquad\qquad \text{if } \beta < \alpha.
\]

\printedpage{187}
% head: A. Grothendieck : Dérivateurs, Ch. XVIII, Édition des Universités de Montpellier II et Paris VII
% foot: Version 17 décembre 2025

\msp{page 483}
Thus, the transfinite induction continues over all of $I$. We see by transfinite induction that we have, for every $\alpha \in I$,
\[
s_\alpha \ \leqslant\ x_\alpha
\]
(card $x_\alpha > x_\beta \geqslant s_\beta$ for $\beta < \alpha$, hence $x_\alpha > s_\beta$ for every $\beta < \alpha$, which is equivalent to $x_\alpha \geqslant s_\alpha$). Moreover, we also find, by transfinite induction, that
\[
\beta \ <\ \alpha \qquad\quad \Longrightarrow \qquad\quad
\text{\parbox[c]{5.5cm}{$s_\beta$ belongs to the closed subcategory $\uS_\alpha$ of $\cM$ generated by the $x_{\beta'}$ with $\beta' < \alpha$.}}
\quad ,
\tag*{$(*)$}
\]
since by construction $s_\alpha$ is the \textit{first} of the elements of $\uS$ which is \textit{not} in the said subcategory $\uS_\alpha$ of $\cM$. This shows that every $s_\beta$ belongs to the closed subcategory ${}_{\underline{x}}\cM$ of $\cM$ generated by the family of \textit{all} the $x_\alpha$, so we find that $\uS \subset {}_{\underline{x}}\cM \subset \uS$, hence $\uS = {}_{\underline{x}}\cM$, q.e.d.

\begingroup\itshape
\noindent\textbf{\upshape Proposition 4.16.7.5.} For a given universe $\U$ (containing $\mathcal{N}$ as an element), the following conditions are equivalent:
\begin{enumerate}[label={}]
\item[1)] There exists a $\U$-accessible category $\cM$ having a closed subcategory not accessible % [sic] 1) -> a), la démonstration renvoie à a) et b)
\msp{page 484}
in $\cM$.
\item[b)] There exists a category $\cC \in \U$, and a transfinite sequence of elements of
\[
\cC^\wedge \ \overset{\text{\upshape def}}{=}\ \underline{\mathrm{Hom}}(\cC, \mathrm{Ens}_{\U}) \ ,
\]
$\underline{x} = (x_\alpha)_{\alpha\in I}$, indexed by the set of all ordinals $\alpha \in \U$, and such that
\[
\Hom(x_\alpha, x_\beta) \ =\ \emptyset \qquad\qquad \textit{if } \alpha < \beta.
\]
\end{enumerate}
\endgroup

\textsc{Proof}. We have b) $\Longrightarrow$ a) by taking $\cM = \cC^\wedge$, by virtue of the preceding proposition. We see likewise a) $\Longrightarrow$ b), by embedding $\cM$ in a category $\cC^\wedge$ ($\cC$ being a small strictly generating subcategory of $\cM$).

I have not been able (\footnote{But yes, in the end, cf. 4.16.7.10 below.}) to construct a sequence $(x_\alpha)_{\alpha\in I}$ as required in this proposition, nor to prove that such a sequence does not exist, even in a particular case such as $\cM = (k\text{-comm. alg.})$, $k$ a commutative ring. I have not
\msp{page 485}
even proved that one cannot find a transfinite sequence of \textit{fields} $k_\alpha$, having the stated properties.

A fortiori, I have not been able to clear up the question whether an $\underline{\mathrm{L}}$-stable subcategory of an accessible category $\cM$ is always accessible in $\cM$ (\footnote{This is false, cf. below.}). It could be that there is no such case. It could also be that there are some, even with a closed subcategory (which is stable under every type of $\varinjlim$ representable in $\cM$, except for the $\varinjlim$ over an empty index category $\ldots$).

\printedpage{188}
% head: A. Grothendieck : Dérivateurs, Ch. XVIII, Édition des Universités de Montpellier II et Paris VII
% foot: Version 17 décembre 2025

For \textit{open} subcategories, I recall that I constructed (in the appendix) an example of such a subcategory, not accessible in $\cM$. Concerning open subcategories, I point out here the statement (almost tautological):

\msp{page 486}
\textbf{Proposition 4.16.7.6} (\footnote{This proposition and its corollary are contained in 4.16.7.8 below, there is no need to make a separate statement of them.}). \textit{Let $\cM$ be an accessible category, $\uS$ an \textup{open} subcategory of $\cM$. For $\uS$ to be accessible in $\cM$, it is necessary and sufficient that it be $\underline{L}$-stable in $\cM$, i.e. that there exist a cardinal $\pi_0$ such that $\cM$ satisfies $L_{\pi_0}$, and that $\uS$ be stable in $\cM$ under $\varinjlim$ of type $L_{\pi_0}$. More precisely, if this condition is satisfied, then for every cardinal $\pi$ such that $\pi \geqslant \pi_0$ and that $\pi$ is useful for $\cM$, $\uS$ is $\pi$-accessible in $\cM$. Or again: if $\pi$ is a useful cardinal for $\cM$, then $\uS$ is $\pi$-accessible in $\cM$ if and only if it is $L_\pi$-stable in $\cM$.}

This last assertion follows at once from the following corollary, itself tautological.

\textbf{Corollary 4.16.7.7.} \textit{Let $\pi$ be a cardinal as indicated above, let $\uS(\pi) = \uS \cap \cM_\pi$, and let $x \in \Ob \cM$ and
\[
x \ \simeq\ \varinjlim_{i\in I} x_i \qquad\qquad
\begin{cases}
x_i \in \Ob \cM_\pi \\[1ex]
I \ \textit{large relative to}\ \pi
\end{cases}
\]
be a representation of $x$ as the essential image of an object of $\Ind_\pi \cM_\pi$. Then $x \in \Ob \uS$ $\Longleftrightarrow$ $x_i \in \Ob \uS(\pi)$ $\forall\, i \in I$.}

\msp{page 487}
I have finally cleared up the question of the accessibility of open subcategories and closed subcategories of an accessible category:

\begingroup\itshape
\noindent\textbf{\upshape Proposition 4.16.7.8.} Let $\cM$ be an accessible category, $\uS$ a closed subcategory of $\cM$, $\uT$ the open subcategory complementary to $\uS$, $\pi$ an infinite cardinal useful for $\cM$. The following conditions are equivalent:
\begin{enumerate}[label={}]
\item[a)] $\uS$ is $\pi$-accessible in $\cM$.
\item[a$'$)] $\uS$ is the closed subcategory generated by $\uS(\pi) \overset{\text{\upshape def}}{=} \uS \cap \cM_\pi$, i.e. for $x \in \Ob \cM$, we have $x \in \Ob \uS$ if and only if $\exists\, \xi \in \Ob \uS(\pi)$ such that $\Hom_{\cM}(\xi, x) \neq \emptyset$.
\item[b)] $\uT$ is $\pi$-accessible in $\cM$.
\item[b$'$)] $\uT$ is $\underline{L}_\pi$-stable in $\cM$.
\end{enumerate}
\endgroup

\noindent Let us note at once the

\begingroup\itshape
\noindent\textbf{\upshape Corollary 4.16.7.9.} For $\cM$, $\uS$, $\uT$ as above, equivalent conditions:
\begin{enumerate}[label={}]
\item[a)] $\uS$ is accessible in $\cM$.
\item[a$'$)] There exists a \textup{small} subset $A$ of $\Ob \cM$, such that $\uS$ is the closed subcategory ${}_A\cM$ generated by $A$, i.e. $\uS = \bigcup_{\xi\in A} {}_\xi\cM$, where ${}_\xi\cM$ is the closed subcategory of $\cM$ consisting of the $x \in \Ob \cM$ such that $\Hom(\xi, x) \neq \emptyset$.
\item[b)] \msp{page 488}$\uT$ is accessible in $\cM$.
\end{enumerate}
\endgroup

\end{document}
